\section{Functional integral and topological sectors}

In lattice gauge theory, the space of gauge fields is connected
and the concept of a topological sector has therefore no
a priori well-defined meaning.
However, one also knows since a long time that any classical
continuous gauge field (including multi-instanton configurations) can
be approximated arbitrarily well by lattice fields.
The sectors are hence included in the field space
but are not separated
from one another.

An understanding of
how exactly the sectors
get divided when the lattice spacing is taken to zero
can now be achieved using the Wilson flow.
The existence of the topological sectors thus
turns out to be a dynamical property of the theory rather than
being a consequence of an assumed or imposed continuity
of the fields.


\subsection{Field transformation}

On a finite lattice, however large, the transformation
\begin{equation}
  U\to V=V_{t_0}
  \label{eq:fieldtrafo}
\end{equation}
is invertible and actually a diffeomorphism of the field space~\cite{TrivMaps}.
One can therefore perform a change of integration
variables in the QCD functional integral from the fundamental field
to the field at flow time $t_0$.
As already noted in ref.~\cite{TrivMaps},
the associated Jacobian can be worked out analytically and be
expressed through the Wilson action along the flow.
The expectation value of any observable $\Obs(V)$ is then
given by
\begin{align}
  \langle\Obs\rangle&={1\over{\cal Z}}\int
  \rmD[V]\,\Obs(V)\,\rme^{-\tilde{S}(V)},
  \label{eq:fint}\\
  \tilde{S}(V)&=S(U)+{16g_0^2\over3a^2}\int_0^{t_0}\rmd t\,\Sw(V_t),
  \qquad
  \rmD[V]=\prod_{x,\mu}\rmd V(x,\mu),
  \label{eq:tSaction}
\end{align}
where $S(U)$ denotes the total action
(including the quark determinants)
of the theory before the transformation.

Evidently, the functional integral may be rewritten in this way
as an integral over the gauge field at any flow time $t$.
Setting $t$ to the reference time $t_0$ is however an
interesting and natural choice.
In particular, the fields that dominate
the transformed integral then have a characteristic
wavelength on the order of the fundamental
low-energy scales of the theory.


\subsection{Smoothness of the dominant fields}

A quantitative measure for the smoothness of a given lattice gauge
field $V$ is
\begin{equation}
  h=\max_{p}s_p,\qquad s_p=\Re\tr\{1-V(p)\},
  \label{eq:hmax}
\end{equation}
the maximum being taken over all plaquettes $p$
(as before, $V(p)$ denotes the product of the link variables
around $p$).
In the functional integral~\eqref{eq:fint}, all possible values of $h$ occur,
but large values are expected to be unlikely
in view of the fact that the expectation value $\langle s_p\rangle$
scales proportionally to $a^4$.

\begin{figure}[ht]
\centering
\includegraphics[width=0.6\textwidth]{images/prob}
\caption{Plot of the probability for
$s_p$ [eq.~\eqref{eq:hmax}] on a given plaquette $p$ to be above
some specified value $s$.
From top to bottom, the three curves shown correspond to
the values $a=0.1$, $0.07$ and $0.05$ fm
of the lattice spacing.
The calculation was performed in the
$\SUthree$ gauge theory using the ensembles
of fields described in subsect.~3.1.}
\label{fig:prob}
\end{figure}

Since all terms in the action~\eqref{eq:tSaction}
have the same sign and favour smooth configurations,
it is certainly plausible that
large plaquette values $s_p$ are strongly suppressed.
The probability for $s_p$ on a given plaquette $p$ to be larger than
some specified value $s$ in fact decreases roughly like $a^{10}$
when the lattice spacing is reduced (see fig.~\ref{fig:prob}).
In a finite volume of fixed physical size,
the statistical weight of the field configurations with
values of $h$ beyond a
given threshold is therefore rapidly going to zero
in the continuum limit, while the fields that dominate
the transformed functional integral
become uniformly smooth in space.


\subsection{Dynamical separation of the topological sectors}

Many years ago, the space of
lattice gauge fields was shown to
divide into disconnected sectors
if only fields satisfying a certain smoothness
condition are admitted~\cite{TopCharge,PhillipsStone}.
The proof is based on a geometrical construction
of a local lattice expression
for the topological charge which assumes integer values and
which has the correct classical continuum limit.

In the case of QCD, the fields satisfying%
\footnote{The theorems proved by Phillips and Stone~\cite{PhillipsStone}
assume a simplicial lattice. To be able to apply them
in the present context, the hypercubic cells of the
lattice must be divided
into simplices. An interpolation of the gauge field to the
added links is then required and can easily be accomplished
using the same interpolation methods as the ones applied elsewhere
in refs.~\cite{TopCharge,PhillipsStone}.}
\begin{equation}
   h<0.067
   \label{eq:hbound}
\end{equation}
are included in the subspace covered by the geometrical construction
and thus fall into topological sectors very much like
the continuous gauge fields in the continuum theory. The space of fields
``between the sectors'' may accordingly be characterized
by the inequality $h\geq 0.067$.

Recalling the discussion in subsect.~4.2
of the smoothness properties of the fields that dominate
the functional integral~\eqref{eq:fint}, the topological sectors
are now seen to emerge as a consequence of the fact that
the fields with large values of $h$ have a rapidly
decreasing weight in the integral when the lattice
spacing is taken zero.
In the case of the simulations reported in sect.~3, for example,
the fraction of representative fields satisfying the bound~\eqref{eq:hbound}
increases from $0\%$ at $a=0.1$ fm to $8\%$ and $70\%$ at
$a=0.07$ and $0.05$ fm, respectively.
The asymptotic behaviour of the weight of the fields
between the sectors cannot be safely determined from these
data, but the curves shown in fig.~\ref{fig:prob} suggest that the
probability of finding such a configuration
on lattices of a given physical size
goes to zero proportionally to $a^6$.


\subsection{Topological susceptibility}

Once the functional integral~\eqref{eq:fint} splits
into a sum of integrals over effectively
disconnected sectors,
the assignment of the topological charge $Q$
to the field configurations
in a representative ensemble of fields
becomes unambiguous.
The geometrical construction mentioned before
could be used to compute the charge, but
in view of the universality property
discussed at the end of sect.~3,
a straightforward discretization of
the topological density, using the symmetric lattice
expression for the field tensor $G_{\mu\nu}$,
is expected to give the same results for the moments
$\langle Q^n\rangle$ in the continuum limit.

On the lattices listed in table~\ref{tab:params},
the second moment $\langle Q^2\rangle$
computed along these lines is about $50$
and the topological susceptibility $\chi_t$ turns out
to be independent of the lattice spacing
within statistical errors. A fit by a constant
then gives the result
\begin{equation}
   \chi_t^{1/4}=187.4(3.9)\,\MeV,
   \label{eq:chit}
\end{equation}
which differs
by less than two standard deviations
from the value $194.5(2.4)$ MeV~\cite{ChiGW}
obtained at flow time $t=0$
using a chiral lattice Dirac operator and
the index theorem~\cite{IndThm}.
Moreover, as shown in a forthcoming publication,
the ratios of the fourth root of the second moments computed here
and through a variant~\cite{ChiProj} of
the universal formula proposed in~\cite{ChiInd} extrapolate to $1$ in the continuum limit
to within a statistical uncertainty of $3\%$.

All these empirical results support the conjecture
that the moments of
the charge distribution in the functional integral~\eqref{eq:fint}
coincide with those obtained from the index theorem~\cite{IndThm,ChiGW,ChiInd,ChiProj}
and thus the
ones appearing in
the chiral Ward identities~\cite{ChiWI,ChiWII}.
However, while highly plausible, the equality remains to be
theoretically established.

